\section{Condensed spectra and solid spectra}\label{sec: condensed spectra and solid spectra}
Lurie's theory of presentable categories produces, from a small site $\Cscr$ and a presentable category $\Vscr$, the hypercomplete $\Vscr$-valued sheaves on $\Cscr$ as the Lurie tensor product of the hypercomplete topos of $\Cscr$ with $\Vscr$. The category of condensed animae is such a hypercomplete topos, and furnishes the basic building block of homotopy-coherent condensed mathematics, and of condensed spectra in particular.

We recall its construction after Clausen \cite[\S 4]{ClausenLie}, building on earlier work of Barwick--Glasman--Haine \cite[\S 13]{Exodromy} and Wolf \cite{WolfPyknotic}, before forming the stable category of condensed spectra by tensoring with $\Sp$. Solid spectra are then isolated within condensed spectra as the analogue of the derived category of solid condensed Abelian groups, and the descriptions of Clausen \cite[Def. 13.2]{ClausenLie} and Wagner \cite[\S B.6]{WagnerHHab} are shown to agree.
\subsection{Condensed animae}\label{subsec: condensed animae}
It was observed by Johnstone that the category $\CHaus$ of compact Hausdorff spaces is monadic over $\Sets$ through the ultrafilter monad \cite[\S III.2.4]{JohnstoneStone}, indicating that a better-behaved topological algebra should follow from the algebraic character of such spaces. This expectation is realized in part by the condensed mathematics of Clausen--Scholze \cite{Condensed}, built from set-valued sheaves on profinite sets for the Grothendieck pretopology of finite jointly surjective covers. 

For most purposes it suffices to restrict to those profinite sets arising as sequential limits of finite discrete spaces, which form an essentially small category and yield the light condensed mathematics used throughout this work. These profinite sets organize into a Grothendieck site whose hypercomplete sheaves are the (light) condensed animae. 
\begin{definition}\label{def: light profinite sets}
    The category of \emph{light profinite sets} $\PFin_{\NN}$ is the sequential $\Pro$-category of finite sets. 
\end{definition}
\begin{lemma}\label{lem: Grothendieck topology on light profinite sets}
    Let $S\in\PFin_{\NN}$ and say a family $\{S_{i}\to S\}_{1\leq i\leq n}$ in $(\PFin_{\NN})_{/S}$ is a covering family if $\coprod_{1\leq i\leq n}S_{i}\to S$ is a surjection. The collection of covering families defines a Grothendieck pretopology on $\PFin_{\NN}$. We refer to the corresponding Grothendieck topology as the finitary Grothendieck topology. 
\end{lemma}
\begin{proof}
    This follows from the discussion in \cite[\S 2.2]{CamargoSolid}. More explicitly, the covering sieves for the finitary Grothendieck topology on $\PFin_{\NN}$ are those sieves that contain a sieve generated by a covering family. 
\end{proof}
\begin{definition}\label{def: condensed animae}
    The category of \emph{condensed animae} $\Cond\Ani$ is the category of hypercomplete $\Ani$-valued sheaves on $\PFin_{\NN}$ with respect to the finitary Grothendieck topology of \Cref{lem: Grothendieck topology on light profinite sets}. 
\end{definition}
Two properties of condensed animae recur throughout. 
\begin{proposition}\label{prop: properties of condensed animae}
    Let $\Cond\Ani$ be the category of condensed animae. 
    \begin{enumerate}
        \item $\Cond\Ani$ is an $\omega_{1}$-presentable topos. 
        \item $\Cond\Ani$ is a replete topos: if 
        $$X:=\lim_{n}\left(\dots\to X_{2}\to X_{1}\to X_{0}\right)$$
        is a sequential diagram in $\Cond\Ani$ with transition maps that are effective epimorphisms then each $X\to X_{n}$ is an effective epimorphism. 
    \end{enumerate}
\end{proposition}
\begin{proof}[Proof of (1)]
    That $\Cond\Ani$ is a topos is immediate from construction. For $\omega_{1}$-presentability, note that the presheaf category $\PSh(\PFin_{\NN};\Ani)$ is $\omega$-presentable so by \cite[Lem. 5.5.1.4]{HTT} it suffices to show the inclusion $\Cond\Ani\hookrightarrow\PSh(\PFin_{\NN};\Ani)$, right adjoint to the formation of hypersheaves, preserves $\omega_{1}$-filtered colimits. Noting that $\omega_{1}$-filtered colimits commute with countable limits in $\Ani$ \cite[Prop. 5.3.3.3]{HTT}, the inclusion preserves them, and the claim follows. 
\end{proof}
\begin{proof}[Proof of (2)]
    This is \cite[Lem. 4.9]{ClausenLie}. 
\end{proof}
\subsection{Condensed spectra}\label{subsec: condensed spectra}
Condensed animae form a presentable topos, and it is by tensoring with spectra that the stable homotopy theory underlying the rest of this work is obtained. The construction rests on the following general identification. 
\begin{proposition}\label{prop: hypersheaves as tensor}
    Let $\Cscr$ be a site with associated hypercomplete topos $\Xscr$ and $\Vscr$ a presentable category. Then precomposition with the functor $\Cscr\to\Xscr$ induces an equivalence of categories $\Xscr\otimes\Vscr\xrightarrow{\sim}\widehat{\Sh}(\Cscr;\Vscr)$. 
\end{proposition}
\begin{proof}
    This follows by combining \cite[Rmk. 1.3.1.6]{SAG}, which identifies $\Xscr\otimes\Vscr$ with $\Vscr$-valued sheaves on $\Xscr$, with \cite[Lem. A.4.22]{HeyerMann6FF}. 
\end{proof}
The derived category of condensed Abelian groups occupies the central position in animated analytic geometry, and condensed spectra take over that role here as its spectral, homotopy-coherent counterpart. 
\begin{definition}\label{def: condensed spectra}
    The category of \emph{condensed spectra} $\Cond\Sp$ is the category of $\Sp$-valued hypersheaves on $\PFin_{\NN}$ with respect to the finitary Grothendieck topology of \Cref{lem: Grothendieck topology on light profinite sets}. 
\end{definition}
\begin{remark}\label{rmk: characterizations of CondSp}
    By \Cref{prop: hypersheaves as tensor} and \cite[Rmk. 1.3.2.2]{SAG}, the category of condensed spectra is defined equivalently as $\Cond\Ani\otimes\Sp$ or as spectrum objects in condensed animae. 
\end{remark}
For the following, recall the notions of an $\omega_{1}$-presentably symmetric monoidal category from \Cref{subsec: notation and conventions} and of a compact projective object from \Cref{def: compact projective objects} (1). With this terminology, the elementary properties of condensed spectra are as follows. 
\begin{theorem}\label{thm: properties of condensed spectra}
    Let $\Cond\Sp$ be the category of condensed spectra. 
    \begin{enumerate}
        \item $\Cond\Sp$ is a stable $\omega_{1}$-presentably symmetric monoidal category with compact projective unit $\Sph$. 
        \item $\Cond\Sp$ admits an accessible $t$-structure $(\Cond\Sp_{\geq0},\Cond\Sp_{\leq0})$ with the following properties: 
        \begin{enumerate}[label=(\roman*)]
            \item $\Cond\Sp_{\geq0},\Cond\Sp_{\leq0}$ are spanned by those condensed spectra $M$ such that the homotopy sheaves $\pi_{k}(M)$ vanish for $k<0$ and $k>0$, respectively. The heart of this $t$-structure is equivalent to the category of (static) condensed Abelian groups $\Cond\Ab$. 
            \item $\Cond\Sp_{\geq0}$ contains the unit in $\Cond\Sp$ and is closed under tensor products in $\Cond\Sp$. 
            \item $\Cond\Sp_{\leq0}$ is closed under filtered colimits in $\Cond\Sp$. 
            \item $\Cond\Sp_{\geq0}$ is closed under countable products in $\Cond\Sp$. 
            \item The $t$-structure on $\Cond\Sp$ is both left and right complete. 
        \end{enumerate}
        \item There is a $t$-exact symmetric monoidal equivalence of categories 
        $$\Mod_{\pi_{0}(\Sph)}(\Cond\Sp)\simeq\Dscr(\Cond\Ab).$$ 
    \end{enumerate}
\end{theorem}
\begin{proof}[Proof of (1)]
    Defining $\Cond\Sp$ as the Lurie tensor product $\Cond\Ani\otimes\Sp$, the statement on $\omega_{1}$-presentability of the underlying category follows readily from the fact that the Lurie tensor product on presentable categories restricts to one on $\PrL_{\kappa}$ for all regular cardinals $\kappa$ \cite[Lem. 5.3.2.11]{HA}. Stability follows from the identification $\Cond\Sp\simeq\Sp(\Cond\Ani)$ in \Cref{rmk: characterizations of CondSp}. The category admits a symmetric monoidal structure by \cite[\S 1.3.4]{SAG}.
    
    For compactness of the unit, we note that every hypercover of the singleton splits, so $\Hom_{\Cond\Sp}(\Sph,-):\Cond\Sp\to\Sp$ preserves filtered colimits which, in conjunction with the finite colimit preservation furnished by exactness, shows the functor preserves small colimits and thus gives compactness of $\Sph$. For projectivity, we show that $\Hom_{\Cond\Sp}(\Sph,-):\Cond\Sp\to\Sp$ given on objects by $M\mapsto M(*)$ is $t$-exact. For $k\in\ZZ$, $\pi_{k}(M(*))\simeq\pi_{k}(M)(*)$ as sheafification of the pointwise homotopy presheaf does not change values at the singleton. Since the $t$-structure on $\Cond\Sp$ is characterized by homotopy sheaves, the functor is $t$-exact. 

    The unit is compact, in particular $\omega_{1}$-compact, so it remains to show that the tensor product preserves $\omega_{1}$-compact objects. The tensor product on $\Cond\Sp$ is exact and preserves small colimits separately in each argument, so using symmetric monoidality of $\Sigma^{\infty}_{+}:\Cond\Ani\to\Cond\Sp$, we have
    $$\Sph[S_{1}]\otimes_{\Cond\Sp}\Sph[S_{2}]\simeq\Sigma^{\infty}_{+}(S_{1})\otimes_{\Cond\Sp}\Sigma^{\infty}_{+}(S_{2})\simeq\Sigma^{\infty}_{+}(S_{1}\times S_{2})\simeq\Sph[S_{1}\times S_{2}]$$
    for $S_{1},S_{2}\in\PFin_{\NN}$. $\PFin_{\NN}$ admits countable limits \cite[Prop. 2.1.4]{CamargoSolid}, so $S_{1}\times S_{2}$ is light profinite and $\Sph[S_{1}\times S_{2}]$ is $\omega_{1}$-compact. The free condensed spectra $\Sph[S]$ for $S\in\PFin_{\NN}$ are $\omega_{1}$-compact and generate $\Cond\Sp$ under small colimits and shifts. Hence an object of $\Cond\Sp$ is $\omega_{1}$-compact if and only if it is contained in the minimal full subcategory of $\Cond\Sp$ closed under countable colimits and shifts and containing the free condensed spectra on light profinite sets -- such a category is already idempotent complete by \cite[Cor. 4.4.5.15]{HTT}. Since the tensor product is exact and preserves colimits separately in each variable, the computation on generators above shows that binary tensors of $\omega_{1}$-compact objects are $\omega_{1}$-compact. Thus $\Cond\Sp$ is $\omega_{1}$-presentably symmetric monoidal. 
\end{proof}
\begin{proof}[Proof of (2)]
    \cite[Prop. 1.3.2.7]{SAG} gives all properties but the condition on countable products and left completeness. 

    \cite[Lem. 4.10 (4)]{ClausenLie} implies that countable products in $\Cond\Sp$ are $t$-exact. In particular, they preserve connective objects, so $\Cond\Sp_{\geq0}$ is closed under countable products. 

    For left completeness, note that the category of condensed animae is a Postnikov complete topos \cite[Thm. A]{PostnikovReplete} as its underlying 1-topos of condensed sets is replete \cite[Prop. 2.1.4 \& Rmk. 2.2.2]{CamargoSolid}. The $t$-structure on $\Sp$-valued sheaves is then left complete by \cite[Cor. 3.26]{PostnikovReplete}. 
\end{proof}
\begin{proof}[Proof of (3)]
    This is \cite[Cor. 2.6.7 \& 2.6.9]{BrinkCondensed}. 
\end{proof}
\begin{remark}\label{rmk: CondSp is pres sym mon t plus bistable}
    The conditions of \Cref{thm: properties of condensed spectra} (1) and (2) show that $\Cond\Sp$ is a presentably symmetric monoidal $t^{+}$-bistable category in the sense of \Cref{def: symmetric monoidal t-bistable category} (2), which situates condensed spectra within the framework of the $t^{+}$-bistable categories discussed in \Cref{app: categories}. 
\end{remark}
\subsection{Solid spectra}\label{subsec: solid spectra}
A solid condensed Abelian group is one in which nullsequences are summable, a condition encoded by the shift operator on the free condensed Abelian group on the nullsequence. The same construction applies to condensed spectra: a condensed spectrum is solid when this shift condition holds for $\Sph[\NN_{\infty}]$, following Wagner \cite[\S B.6]{WagnerHHab}. This coincides with the formulation of Clausen \cite[Def. 13.2]{ClausenLie}, under which solidity is tested on homotopy groups.
\begin{definition}\label{def: solid spectra}
    Let $\NN\cup\{\infty\}$ be the one-point compactification of the natural numbers as a light profinite set and $\Sph[\NN_{\infty}]:=\cofib(\Sph[\{\infty\}]\to\Sph[\NN\cup\{\infty\}])$ be the cofiber of the morphism of condensed spectra induced by the inclusion $\{\infty\}\hookrightarrow\NN\cup\{\infty\}$ in $\PFin_{\NN}\subseteq\Cond\Ani$. 
    \begin{enumerate}
        \item A condensed spectrum $M$ is \emph{solid} if the natural morphism 
        $$\id-\sigma^{*}:\intHom_{\Cond\Sp}(\Sph[\NN_{\infty}],M)\to\intHom_{\Cond\Sp}(\Sph[\NN_{\infty}],M)$$
        induced by the shift operator $\sigma$ given by $n\mapsto n+1$ on $\NN\cup\{\infty\}$ is an equivalence. 
        \item The category of \emph{solid spectra} $\Cond\Sp^{\sld}$ is the full subcategory of $\Cond\Sp$ spanned by the solid spectra. 
    \end{enumerate}
\end{definition}
The following records the colimit behaviour of the internal hom out of the nullsequence quotient, reducing it to Brink's computation for internal cohomology of free condensed spectra on light profinite sets \cite[Lem. 2.7.1]{BrinkCondensed}.
\begin{lemma}\label{lem: Brink internal cohomology}
    Let $\Sph[\NN_{\infty}]:=\cofib(\Sph[\{\infty\}]\to\Sph[\NN\cup\{\infty\}])$ be the object of \Cref{def: solid spectra}. Then the functor $\intHom_{\Cond\Sp}(\Sph[\NN_{\infty}],-):\Cond\Sp\to\Cond\Sp$ commutes with direct sums of uniformly bounded above objects: 
    $$\intHom_{\Cond\Sp}\left(\Sph[\NN_{\infty}],\bigoplus_{i\in I}M_{i}\right)\simeq\bigoplus_{i\in I}\intHom_{\Cond\Sp}\left(\Sph[\NN_{\infty}],M_{i}\right)$$
    for every $k\geq0$ and every family $\{M_{i}\}_{i\in I}\subseteq\Cond\Sp_{\leq k}$. 
\end{lemma}
\begin{proof}
    Applying $\intHom_{\Cond\Sp}(-,-)$ to the (co)fiber sequence $\Sph[\{\infty\}]\to\Sph[\NN\cup\{\infty\}]\to\Sph[\NN_{\infty}]$ in the first argument gives a (co)fiber sequence 
    $$\intHom_{\Cond\Sp}(\Sph[\NN_{\infty}],-)\to\intHom_{\Cond\Sp}(\Sph[\NN\cup\{\infty\}],-)\to\intHom_{\Cond\Sp}(\Sph[\{\infty\}],-)$$
    of endofunctors of $\Cond\Sp$. By \cite[Lem. 2.7.1]{BrinkCondensed}, the latter two functors commute with filtered colimits of uniformly bounded above objects. Since limits and colimits in functor categories are computed pointwise, and finite limits commute with filtered colimits in $\Cond\Sp$, $\intHom_{\Cond\Sp}(\Sph[\NN_{\infty}],-)$ preserves filtered colimits of uniformly bounded above objects. 

    Now in the case of direct sums, we may write $\bigoplus_{i\in I}M_{i}$ as a filtered colimit of its finite subsums. Since $\intHom_{\Cond\Sp}(\Sph[\NN_{\infty}],-)$ is exact it preserves finite direct sums, and additionally preserves filtered colimits, giving the claim. 
\end{proof}
\begin{proposition}\label{prop: nullsequence quotient is internally compact projective}
    Let $\Sph[\NN_{\infty}]:=\cofib(\Sph[\{\infty\}]\to\Sph[\NN\cup\{\infty\}])$ be the object of \Cref{def: solid spectra}. The object $\Sph[\NN_{\infty}]$ is an internally compact projective object of $\Cond\Sp$ in the sense of \Cref{def: internally compact projective}: $\intHom_{\Cond\Sp}(\Sph[\NN_{\infty}],-):\Cond\Sp\to\Cond\Sp$ is $t$-exact and preserves small colimits. 
\end{proposition}
\begin{proof}
    We first treat $t$-exactness. Note $\Sph[\NN_{\infty}]$ is connective, being the cofiber of connective objects, so $\intHom_{\Cond\Sp}(\Sph[\NN_{\infty}],-)$ is exact, limit-preserving, and left $t$-exact as an endofunctor of $\Cond\Sp$ as it is right adjoint to the right $t$-exact left adjoint $\Sph[\NN_{\infty}]\otimes_{\Cond\Sp}-$ \cite[Prop. 1.3.17 (ii)]{Perverse}. The equivalence of \Cref{thm: properties of condensed spectra} (3), being $t$-exact, identifies the respective hearts. Under this identification, \Cref{lem: modules inherit bistability} (2) and \Cref{prop: relationship with heart} (1) identify the restriction of $\intHom_{\Cond\Sp}(\Sph[\NN_{\infty}],-)$ to the heart with $\intHom_{\Dscr(\Cond\Ab)}(\ZZ[\NN_{\infty}],-)|_{\Sph}$. By internal projectivity of $\ZZ[\NN_{\infty}]$ in $\Cond\Ab$ \cite[Thm. 2.3.3]{CamargoSolid}, the latter agrees with $\intHom_{\Cond\Ab}(\ZZ[\NN_{\infty}],-)|_{\Sph}$ \cite[Rmk. 3.2.2]{CamargoSolid} and hence preserves the heart. Furthermore, $\Cond\Sp$ is $\lim^{1}$-epi-acyclic by \Cref{rmk: 1-localic is lim1 epi acyclic}, its heart being Abelian sheaves on the replete 1-topos of condensed sets \cite[Prop. 2.1.4 \& Rmk. 2.2.2]{CamargoSolid}, so the hypotheses of \Cref{prop: lim1 acyclic t-exact} are satisfied and $\intHom_{\Cond\Sp}(\Sph[\NN_{\infty}],-)$ is $t$-exact.

    By exactness of $\intHom_{\Cond\Sp}(\Sph[\NN_{\infty}],-)$, preservation of small colimits can be reduced to the preservation of direct sums. Note that $\Cond\Sp_{\leq k}$ is closed under arbitrary direct sums in $\Cond\Sp$. Indeed, every direct sum is the filtered colimit of its finite subsums, the latter belong to $\Cond\Sp_{\leq k}$, and $\Cond\Sp_{\leq k}$ is closed under filtered colimits by compatibility with filtered colimits \Cref{thm: properties of condensed spectra} (2) (iii). Hence direct sums of objects of $\Cond\Sp_{\leq k}$ agree with those in $\Cond\Sp$. For a direct sum $\bigoplus_{i\in I}M_{i}$, we have $\tau_{\leq k}\bigoplus_{i\in I}M_{i}\simeq\bigoplus_{i\in I}\tau_{\leq k}M_{i}$ by commuting Postnikov truncation, a left adjoint, with the colimit so $\bigoplus_{i\in I}M_{i}\simeq\lim_{k}\tau_{\leq k}\bigoplus_{i\in I}M_{i}\simeq\lim_{k}\bigoplus_{i\in I}\tau_{\leq k}M_{i}$ by left completeness. We then compute 
    \begin{align*}
        \intHom_{\Cond\Sp}\left(\Sph[\NN_{\infty}],\bigoplus_{i\in I}M_{i}\right) &\simeq \lim_{k}\intHom_{\Cond\Sp}\left(\Sph[\NN_{\infty}],\bigoplus_{i\in I}\tau_{\leq k}M_{i}\right) \\
        &\simeq\lim_{k}\bigoplus_{i\in I}\intHom_{\Cond\Sp}\left(\Sph[\NN_{\infty}],\tau_{\leq k}M_{i}\right) && \text{\Cref{lem: Brink internal cohomology}} \\
        &\simeq \lim_{k}\bigoplus_{i\in I}\tau_{\leq k}\intHom_{\Cond\Sp}\left(\Sph[\NN_{\infty}],M_{i}\right) && t\text{-exactness} \\
        &\simeq \lim_{k}\tau_{\leq k}\bigoplus_{i\in I}\intHom_{\Cond\Sp}\left(\Sph[\NN_{\infty}],M_{i}\right) && \tau_{\leq k}\text{ left adj.} \\
        &\simeq \bigoplus_{i\in I}\intHom_{\Cond\Sp}(\Sph[\NN_{\infty}],M_{i}) && \text{left compl.}
    \end{align*}
    yielding the claim. 
\end{proof}
With this key computation in hand, the closure properties and the homotopy group criterion follow directly.
\begin{theorem}\label{thm: properties of solid spectra}
    Let $\Cond\Sp^{\sld}$ be the category of solid spectra. 
    \begin{enumerate}
        \item $\Cond\Sp^{\sld}$ is closed under small limits, small colimits, and extensions in $\Cond\Sp$. 
        \item If $M\in\Cond\Sp$ then $M\in\Cond\Sp^{\sld}$ if and only if $\pi_{k}(M)\in\Cond\Ab^{\sld}$ for all $k\in\ZZ$. 
    \end{enumerate}
\end{theorem}
\begin{proof}[Proof of (1)]
    The functor $\intHom_{\Cond\Sp}(\Sph[\NN_{\infty}],-)$ is a right adjoint, hence preserves small limits. Closure of $\Cond\Sp^{\sld}$ under small colimits follows from internal compact projectivity shown in \Cref{prop: nullsequence quotient is internally compact projective}. For closure under extensions, we let $M_{1}\to M_{2}\to M_{3}$ be a (co)fiber sequence in $\Cond\Sp$ with $M_{1},M_{3}\in\Cond\Sp^{\sld}$. Contemplate the diagram 
    $$% https://q.uiver.app/#q=WzAsNixbMCwwLCJNX3sxfSJdLFswLDEsIjAiXSxbMiwwLCJNX3syfSJdLFsyLDEsIk1fezN9Il0sWzQsMSwiTV97MX1bMV0iXSxbNCwwLCIwIl0sWzAsMl0sWzIsNV0sWzUsNF0sWzIsM10sWzMsNF0sWzAsMV0sWzEsM11d
    \begin{tikzcd}
        {M_{1}} && {M_{2}} && 0 \\
        0 && {M_{3}} && {M_{1}[1]}
        \arrow[from=1-1, to=1-3]
        \arrow[from=1-1, to=2-1]
        \arrow[from=1-3, to=1-5]
        \arrow[from=1-3, to=2-3]
        \arrow[from=1-5, to=2-5]
        \arrow[from=2-1, to=2-3]
        \arrow[from=2-3, to=2-5]
    \end{tikzcd}$$
    where both squares are biCartesian. Since $M_{1}[1],M_{3}\in\Cond\Sp^{\sld}$ so is $M_{2}$, being the fiber of $M_{3}\to M_{1}[1]$. 
\end{proof}
\begin{proof}[Proof of (2)]
    Suppose $M\in\Cond\Sp^{\sld}$. By $t$-exactness of $\intHom_{\Cond\Sp}(\Sph[\NN_{\infty}],-)$ we have 
    $$\pi_{k}\left(\intHom_{\Cond\Sp}(\Sph[\NN_{\infty}],M)\right)\simeq\intHom_{\Cond\Ab}(\ZZ[\NN_{\infty}],\pi_{k}(M))$$
    where the equivalence follows from \Cref{lem: modules inherit bistability} (2) and internal projectivity of $\ZZ[\NN_{\infty}]$ in $\Cond\Ab$, which identifies the derived internal hom with the underived internal hom \cite[Thm. 2.3.3 \& Rmk. 3.2.2]{CamargoSolid}. This equivalence is natural in $\id-\sigma^{*}$, so solidity of $M$ implies that $\pi_{k}(M)\in\Cond\Ab^{\sld}$ for every $k$. The same computation applied to an object $A\in\Cond\Ab\simeq\Cond\Sp^{\heartsuit}$ shows that $A\in\Cond\Ab^{\sld}$ if and only if $A$, regarded as a condensed spectrum, belongs to $\Cond\Sp^{\sld}$.

    Conversely suppose $\pi_{k}(M)\in\Cond\Ab^{\sld}$ for all $k\in \ZZ$. By right completeness and closure of $\Cond\Sp^{\sld}$ under small colimits from (1), it suffices to treat $M$ bounded below, and after shifting we may assume that $M$ is connective. By left completeness, we write $M\simeq\lim_{n}\tau_{\leq n}M$. Since $\Cond\Sp^{\sld}$ is closed under limits, it suffices to show that each term of the Postnikov tower lies in $\Cond\Sp^{\sld}$. For $n\geq 1$ the terms of the Postnikov tower form (co)fiber sequences 
    $$\pi_{n}(M)[n]\to\tau_{\leq n}M\to\tau_{\leq n-1}M$$
    from which the desideratum is obtained by induction: for $n=1$ the (co)fiber sequence is $\pi_{1}(M)[1]\to\tau_{\leq 1}M\to\pi_{0}(M)$ so $\tau_{\leq 1}M\in\Cond\Sp^{\sld}$ by closure under extensions. For larger $n$ we argue similarly, for the leftmost term is in $\Cond\Sp^{\sld}$ by assumption and the rightmost by the induction hypothesis. 
\end{proof}
\begin{remark}\label{rmk: definitions of solid spectra agree}
    \Cref{thm: properties of solid spectra} (2) shows that the definitions of solid spectra given by Wagner in \cite[\S B.6]{WagnerHHab} and Clausen in \cite[Def. 13.2]{ClausenLie} coincide. More generally, this is a manifestation of the fact that sufficiently nice localizations of stable presentable categories are determined by their heart. See \Cref{subsec: localizations of stable categories} for variations on this theme. 
\end{remark}